
\lussec 7. 关联函数的计算

在数值格点 QCD 中，
涉及非零流时间处的场的关联函数，
可以按照处理普通强子关联函数时通常采取的步骤来计算。
然而流方程必须在某个环节被求解，
这一事实代表了一种需要仔细对待的复杂化。
作为示例，本节针对两个代表性情形把细节具体展开。


\lussubsec 7.1 赝标两点函数

实践中人们关心的是零空间动量下的关联函数
\equation{
  \langle \Pax^{rs}(x)\Pax_t^{sr}(y)\rangle=
  -\sum_{v,w}\langle\tr\{
  [K(t,y;0,v)S(v,x)_{ss}]^{\dagger}K(t,y;0,w)S(w,x)_{rr}\}\rangle
  \enum
}
并利用随机源场 [\cite{MichaelPeisa}]
计算其对一个源点格点集合 $\Gamma$ 的平均，
以减小统计误差。

若取 $x$ 为源点，
则求平均相当于作代换
\equation{
  P^{rs}(x)\to {1\over N_{\Gamma}N_s}
  \sum_{k=1}^{N_s}
  (\psibar_r,\eta_k)(\eta_k,\dirac{5}\psi_s),
  \enum
}
其中 $N_{\Gamma}$ 是 $\Gamma$ 中点的数目，
而
\equation{
  \eta_k(x),\quad k=1,\ldots,N_s,
  \enum
}
是 $\Gamma$ 上随机选取的、均值为零且方差为
\equation{
  \langle\eta_k(v)\eta_l(w)^{\dagger}\rangle=
  \cases{\delta_{kl}\delta_{vw} & if $v,w\in\Gamma$,\cr
  \noalign{\vskip1.0ex}
  0 & otherwise.\cr}
  \enum
}
的复旋量场。
在一个 QCD 模拟中，
计算零空间动量下平均后的关联函数
\equation{
  {1\over N_{\Gamma}}\sum_{x\in\Gamma}\sum_{\vec{y}}
  \langle \Pax^{rs}(x)\Pax_t^{sr}(y)\rangle=
  -{1\over N_{\Gamma}N_s}\sum_{k=1}^{N_s}
  \sum_{\vec{y}}\langle\phi_{k,s}(t,y)^{\dagger}\phi_{k,r}(t,y)\rangle,
  \enum
}
于是相当于对有代表性的规范场系综、
对所有 $k=1,\ldots,N_s$ 与 $h=r,s$，
计算函数
\equation{
  \phi_{k,h}(t,y)=\sum_{w,x}
  K(t,y;0,w)S(w,x)_{hh}\eta_k(x)
  \enum
}
对给定的规范场，
计算 $\phi_{k,h}(t,y)$ 分两步进行：
先通过求解质量为 $m_{0,h}$、源场为 $\eta_k(x)$
的格点狄拉克方程算出 $\phi_{k,h}(0,w)$；
随后必须把算得的场在流时间上从 0 演化到 $t$。
沿流时间增大的方向，
流方程在数值上是稳定的，
用龙格--库塔积分器（附录 D）
便容易以可忽略的积分误差得到解。


\lussubsec 7.2 手征凝聚

对随时间演化的夸克凝聚
\equation{
  \langle S_t^{rr}(x)\rangle=
  -\sum_{v,w}\langle
  \tr\{K(t,x;0,v)\left[S(v,w)_{rr}-\cfl\delta_{vw}\right]
  K(t,x;0,w)^{\dagger}\}\rangle,
  \enum
}
而言，同样希望对其位置 $x$ 求平均。
如上引入随机源场后，即可得到表示式
\equation{
  {1\over N_{\Gamma}}\sum_{x\in\Gamma}
  \langle S_t^{rr}(x)\rangle=
  -{1\over N_{\Gamma}}
  \sum_{v,w}\langle
  \xi_k(t;0,v)^{\dagger}\left[S(v,w)_{rr}-\cfl\delta_{vw}\right]
  \xi_k(t;0,w)\rangle
  \enum
}
其中用到场
\equation{
  \xi_k(t;s,w)=\sum_x K(t;x;s,w)^{\dagger}\eta_k(x)
  \enum
}
在流时间 $s=0$ 处的值。

计算这些场需要求解伴随流方程
\equation{
  (\partial_s+\Delta)\xi_k(t;s,w)=0
  \enum
}
从时刻 $s=t$（此时 $\xi_k$ 与 $\eta_k$ 重合）
积到 $s=0$。
此处不应使用直截了当的龙格--库塔积分，
因为拉普拉斯算子 $\Delta$ 是在由 $s$ 时刻梯度流决定的规范场存在下取的，
那样就必须把规范场沿时间反向演化，
即沿不稳定方向演化。
如附录 E 所解释的，
这一困难可以通过一个避免规范场反向积分的分级方案来克服。
